% To merge into the real deck: replace the next line with
% % !TeX root = raytracing/slides/main.tex
% raytracing_preamble.tex
% Common styling for Ray Tracing presentation

\documentclass[10pt]{beamer}
\usetheme{metropolis}
\usefonttheme{professionalfonts}

% --- Color Definitions ---
\definecolor{PrimaryColor}{RGB}{33,52,72}         % Deep navy
\definecolor{SecondaryColor}{RGB}{84,119,146}     % Desaturated blue-gray
\definecolor{AccentColor}{RGB}{100, 156, 165}       % Soft steel blue

\definecolor{BackgroundColor}{RGB}{245,247,250}   % Very light cool gray/blue-tinted white
\definecolor{TextColor}{RGB}{40,55,70}            % Darker, cool-toned charcoal (slightly less saturated)
\definecolor{LightGray}{RGB}{220,225,230}         % Soft, cool light gray
\definecolor{DarkGray}{RGB}{70,90,105}            % Cool-toned dark gray (coherent with Primary/Secondary)

\definecolor{RayColor}{RGB}{230,150,80}           % Muted warm orange (less saturated than pure orange)
\definecolor{ObjectColor}{RGB}{128,101,160}       % Muted violet (adjusted purple to match cool palette)
\definecolor{LightColor}{RGB}{235,200,100}        % Soft golden yellow (to complement cool blues)

% --- Theme Customization ---
\setbeamercolor{background canvas}{bg=BackgroundColor}
\setbeamercolor{normal text}{fg=TextColor}
\setbeamercolor{frametitle}{bg=PrimaryColor, fg=white}
\setbeamercolor{section in toc}{fg=PrimaryColor}
\setbeamercolor{block title}{bg=PrimaryColor!80, fg=white}
\setbeamercolor{block body}{bg=PrimaryColor!10}
\setbeamercolor{alerted text}{fg=AccentColor}
\setbeamercolor{itemize item}{fg=PrimaryColor}
\setbeamercolor{itemize subitem}{fg=SecondaryColor}
\setbeamerfont{frametitle}{size=\large,series=\bfseries}

% --- Packages ---
\usepackage[utf8]{inputenc}
\usepackage{url}
\usepackage{booktabs}
\usepackage{amsmath, amssymb}
\usepackage{fontawesome5}
\usepackage{pifont}
\usepackage[most]{tcolorbox}
\tcbuselibrary{skins}
\usepackage{colortbl}
\usepackage{array}
\usepackage{tikz}
\usepackage{graphicx}
\usetikzlibrary{shapes.callouts, positioning, arrows.meta, shapes.geometric, shadows, calc, patterns, 3d, backgrounds, shadings}
\usepackage{adjustbox}
\usepackage{ragged2e}
\usepackage{pgfplots}
\usepackage{graphicx}
\usepackage{caption}
\usepackage{minted}
\usepackage{subcaption}
\pgfplotsset{compat=1.18}

% --- Custom Commands ---
\newcommand{\cmark}{\textcolor{SecondaryColor}{\ding{51}}}
\newcommand{\xmark}{\textcolor{AccentColor}{\ding{55}}}
\newcommand{\highlight}[1]{\textcolor{PrimaryColor}{\textbf{#1}}}
\newcommand{\raycolor}[1]{\textcolor{RayColor}{\textbf{#1}}}
\newcommand{\objectcolor}[1]{\textcolor{ObjectColor}{\textbf{#1}}}

% --- TikZ styles for ray tracing visualizations ---
\tikzstyle{process} = [rectangle, rounded corners=3mm, minimum width=2cm, minimum height=0.8cm, text centered, draw=PrimaryColor, thick, fill=PrimaryColor!15, drop shadow]
\tikzstyle{arrow} = [thick, PrimaryColor, ->, >=stealth]
\tikzstyle{ray} = [thick, RayColor, ->, >=stealth]
\tikzstyle{lightray} = [thick, LightColor, ->, >=stealth]
\tikzstyle{reflectray} = [thick, SecondaryColor, ->, >=stealth]
\tikzstyle{refractray} = [thick, AccentColor, ->, >=stealth]
\tikzstyle{shadowray} = [thick, DarkGray, ->, >=stealth, dashed]

% --- 3D object styles ---
\tikzstyle{sphere} = [circle, minimum size=1.5cm, draw=ObjectColor, thick, fill=ObjectColor!20, drop shadow]
\tikzstyle{plane} = [rectangle, minimum width=3cm, minimum height=0.2cm, draw=ObjectColor, thick, fill=ObjectColor!20]
\tikzstyle{triangle} = [regular polygon, regular polygon sides=3, minimum size=1.5cm, draw=ObjectColor, thick, fill=ObjectColor!20]

% --- Eye/Camera styles ---
\tikzstyle{eye} = [circle, minimum size=0.8cm, draw=PrimaryColor, thick, fill=PrimaryColor!30]
\tikzstyle{pixel} = [rectangle, minimum size=0.2cm, draw=AccentColor, fill=AccentColor!30]

% --- Text box styles ---
\tikzstyle{conceptbox} = [rectangle, rounded corners, fill=PrimaryColor!10, draw=PrimaryColor, thick, text width=0.8\textwidth, inner sep=8pt]
\tikzstyle{formulabox} = [rectangle, rounded corners, fill=SecondaryColor!10, draw=SecondaryColor, thick, inner sep=6pt]

% --- Custom environments ---
\newtcolorbox{raybox}[1]{
  colback=RayColor!10,
  colframe=RayColor,
  title=#1,
  fonttitle=\bfseries,
  sharp corners
}

\newtcolorbox{conceptbox}[1]{
  colback=PrimaryColor!10,
  colframe=PrimaryColor,
  title=#1,
  fonttitle=\bfseries,
  rounded corners
}

\newtcolorbox{mathbox}[1]{
  colback=SecondaryColor!10,
  colframe=SecondaryColor,
  title=#1,
  fonttitle=\bfseries,
  rounded corners
}

\newenvironment{timeline}{%
  \begin{tikzpicture}[scale=0.8]
    \coordinate (start) at (0,0);
    \newcounter{timelineitem}
    \setcounter{timelineitem}{0}
  }{%
  \end{tikzpicture}
}

\newcommand{\timelineitem}[4]{%
  \stepcounter{timelineitem}%
  % Name of this item’s node
  \edef\thisID{time\thetimelineitem}%
  \ifnum\value{timelineitem}=1
  % First item: anchor at (start)
  \node[rectangle, draw, fill=PrimaryColor!20,
  text width=1.5cm, minimum height=0.8cm, align=center]
  (\thisID) at (start) {\textbf{#2}};
  \else
  % Later items: below=#1 of previous time node
  \pgfmathtruncatemacro\prev{\value{timelineitem}-1}%
  \node[rectangle, draw, fill=PrimaryColor!20,
    text width=1.5cm, minimum height=0.8cm, align=center,
  below=#1 of time\prev]
  (\thisID) {\textbf{#2}};
  \fi
  % Description always to the right of this time node
  \node[rectangle, draw, fill=SecondaryColor!10,
    text width=4cm, minimum height=0.8cm,
  align=left, right=0.3cm of \thisID]
  (desc\thetimelineitem) {\textbf{#3}\\#4};
  % Draw arrow back to previous if not the first
  \ifnum\value{timelineitem}>1
  \draw[->, thick, PrimaryColor]
  (time\prev.south) -- (\thisID.north);
  \fi
}

\tikzset{
  camera/.style={fill=PrimaryColor!60, draw=PrimaryColor!80, rectangle, minimum size=8pt},
  image plane/.style={fill=AccentColor!10, draw=AccentColor!50, opacity=0.8},
  pixel/.style={fill=AccentColor!60, thick},
  primary ray/.style={->, very thick, red!90},
  object/.style={fill=ObjectColor!60, draw=ObjectColor!80, circle, minimum size=12pt},
  fovangle/.style={<->, thick, PrimaryColor, dashed}
}

\tikzset{
  lens/.style={thick, PrimaryColor, line width=3pt},
  focal plane/.style={thick, AccentColor},
  object ray/.style={->, thick, ObjectColor},
  image ray/.style={->, thick, SecondaryColor},
  optical axis/.style={dashed, gray}
}

\newcommand{\lens}[3]{%
  \begingroup
  % half‐dimensions
  \pgfmathsetlengthmacro{\a}{#2}%
  \pgfmathsetlengthmacro{\b}{#3}%
  \pgfmathsetlengthmacro{\c}{#3/2}%
  \begin{scope}[shift={#1}]
    \draw[line join=round, fill=blue!15]
    (0,-{\c})
    arc(-30:30:{\a} and {\b})
    arc(150:210:{\a} and {\b})
    ;
  \end{scope}
  \endgroup
}

% ================================================================
% Reference theme, rebuilt to match combined/deck.pdf
% Rust title bar, cream ground, Adventor headings, thin footer.
% Placed last so it overrides metropolis.
% ================================================================
\usepackage{tgheros}
\usepackage{tgadventor}
\renewcommand{\sfdefault}{qhv}                       % Heros for body copy
\newcommand{\headfont}{\fontfamily{qag}\selectfont}  % Adventor for headings

\definecolor{ThemeRust}{HTML}{9B4632}
\definecolor{ThemeCream}{HTML}{F8F6F2}
\definecolor{ThemeSage}{HTML}{C6CCB1}
\definecolor{ThemeStone}{HTML}{E7E7DB}
\definecolor{ThemeMuted}{HTML}{8B8A88}
\definecolor{ThemeInk}{HTML}{191919}

\setbeamercolor{background canvas}{bg=ThemeCream}
\setbeamercolor{normal text}{fg=ThemeInk}
\setbeamercolor{frametitle}{bg=ThemeRust, fg=ThemeCream}
\setbeamercolor{footer}{bg=ThemeCream, fg=ThemeMuted}
\setbeamercolor{alerted text}{fg=ThemeRust}
\setbeamercolor{itemize item}{fg=ThemeRust}
\setbeamercolor{itemize subitem}{fg=ThemeRust}
\setbeamercolor{palette primary}{bg=ThemeRust, fg=ThemeCream}  % section pages

% full-bleed rust bar across the top
\setbeamertemplate{frametitle}{%
  \nointerlineskip%
  \begin{beamercolorbox}[wd=\paperwidth,sep=0pt]{frametitle}%
    \vbox{%
      \vskip0.60em
      \hbox{\hskip0.9em%
        \begin{minipage}{\dimexpr\paperwidth-1.8em\relax}%
          \raggedright\headfont\bfseries\large\insertframetitle%
        \end{minipage}}%
      \vskip0.55em
    }%
  \end{beamercolorbox}%
}

% thin three-zone footer
\setbeamertemplate{footline}{%
  \hbox{%
  \begin{beamercolorbox}[wd=.40\paperwidth,ht=1.9ex,dp=1.0ex,left]{footer}%
    \hspace*{0.9em}\tiny\insertshortauthor
  \end{beamercolorbox}%
  \begin{beamercolorbox}[wd=.28\paperwidth,ht=1.9ex,dp=1.0ex,center]{footer}%
    \tiny\insertshorttitle
  \end{beamercolorbox}%
  \begin{beamercolorbox}[wd=.32\paperwidth,ht=1.9ex,dp=1.0ex,right]{footer}%
    \tiny\insertframenumber\,/\,\inserttotalframenumber\hspace*{0.9em}
  \end{beamercolorbox}}%
  \vskip1pt
}

% full-bleed rust title page
\setbeamertemplate{title page}{%
  \begin{tikzpicture}[remember picture, overlay]
    \fill[ThemeRust] (current page.south west) rectangle (current page.north east);
  \end{tikzpicture}%
  \vfill
  \begin{flushright}
    {\headfont\bfseries\LARGE\color{ThemeCream}\inserttitle\par}
    \vspace{0.3em}
    {\color{ThemeCream!85}\small\insertsubtitle\par}
    \vspace{1.3em}
    {\color{ThemeCream}\small\insertauthor\par}
    \vspace{0.5em}
    {\color{ThemeCream!80}\scriptsize\insertinstitute\par}
    \vspace{0.8em}
    {\color{ThemeCream!75}\scriptsize\insertdate\par}
  \end{flushright}
  \vfill
}

% --- Deck palette remapped onto the theme -------------------------------
% Overrides the navy/violet set defined at the top of this file. Every TikZ
% figure and tcolorbox references these names, so they all follow along.
\definecolor{PrimaryColor}{HTML}{9B4632}     % rust: highlights, frames, titles
\definecolor{SecondaryColor}{HTML}{7A5C4B}   % warm brown: secondary boxes
\definecolor{AccentColor}{HTML}{7E8C57}      % sage: arrows, alerted text
\definecolor{ObjectColor}{HTML}{A9745B}      % clay: solid shapes
\definecolor{LightColor}{HTML}{D9A441}       % warm gold
\definecolor{TextColor}{HTML}{191919}
\definecolor{BackgroundColor}{HTML}{F8F6F2}
\definecolor{LightGray}{HTML}{E7E7DB}
\definecolor{DarkGray}{HTML}{6B6257}

% the four fern rules -- must match COLS in part 3/figs5.py
\definecolor{RuleA}{HTML}{C0392B}
\definecolor{RuleB}{HTML}{2E7D32}
\definecolor{RuleC}{HTML}{B8860B}
\definecolor{RuleD}{HTML}{1F6F8B}

  (the teacher's actual preamble)
% Compile from the repo root:  pdflatex -shell-escape "part 3/main.tex"
% !TeX root = raytracing/slides/main.tex
% raytracing_preamble.tex
% Common styling for Ray Tracing presentation

\documentclass[10pt]{beamer}
\usetheme{metropolis}
\usefonttheme{professionalfonts}

% --- Color Definitions ---
\definecolor{PrimaryColor}{RGB}{33,52,72}         % Deep navy
\definecolor{SecondaryColor}{RGB}{84,119,146}     % Desaturated blue-gray
\definecolor{AccentColor}{RGB}{100, 156, 165}       % Soft steel blue

\definecolor{BackgroundColor}{RGB}{245,247,250}   % Very light cool gray/blue-tinted white
\definecolor{TextColor}{RGB}{40,55,70}            % Darker, cool-toned charcoal (slightly less saturated)
\definecolor{LightGray}{RGB}{220,225,230}         % Soft, cool light gray
\definecolor{DarkGray}{RGB}{70,90,105}            % Cool-toned dark gray (coherent with Primary/Secondary)

\definecolor{RayColor}{RGB}{230,150,80}           % Muted warm orange (less saturated than pure orange)
\definecolor{ObjectColor}{RGB}{128,101,160}       % Muted violet (adjusted purple to match cool palette)
\definecolor{LightColor}{RGB}{235,200,100}        % Soft golden yellow (to complement cool blues)

% --- Theme Customization ---
\setbeamercolor{background canvas}{bg=BackgroundColor}
\setbeamercolor{normal text}{fg=TextColor}
\setbeamercolor{frametitle}{bg=PrimaryColor, fg=white}
\setbeamercolor{section in toc}{fg=PrimaryColor}
\setbeamercolor{block title}{bg=PrimaryColor!80, fg=white}
\setbeamercolor{block body}{bg=PrimaryColor!10}
\setbeamercolor{alerted text}{fg=AccentColor}
\setbeamercolor{itemize item}{fg=PrimaryColor}
\setbeamercolor{itemize subitem}{fg=SecondaryColor}
\setbeamerfont{frametitle}{size=\large,series=\bfseries}

% --- Packages ---
\usepackage[utf8]{inputenc}
\usepackage{url}
\usepackage{booktabs}
\usepackage{amsmath, amssymb}
\usepackage{fontawesome5}
\usepackage{pifont}
\usepackage[most]{tcolorbox}
\tcbuselibrary{skins}
\usepackage{colortbl}
\usepackage{array}
\usepackage{tikz}
\usepackage{graphicx}
\usetikzlibrary{shapes.callouts, positioning, arrows.meta, shapes.geometric, shadows, calc, patterns, 3d, backgrounds, shadings}
\usepackage{adjustbox}
\usepackage{ragged2e}
\usepackage{pgfplots}
\usepackage{graphicx}
\usepackage{caption}
\usepackage{minted}
\usepackage{subcaption}
\pgfplotsset{compat=1.18}

% --- Custom Commands ---
\newcommand{\cmark}{\textcolor{SecondaryColor}{\ding{51}}}
\newcommand{\xmark}{\textcolor{AccentColor}{\ding{55}}}
\newcommand{\highlight}[1]{\textcolor{PrimaryColor}{\textbf{#1}}}
\newcommand{\raycolor}[1]{\textcolor{RayColor}{\textbf{#1}}}
\newcommand{\objectcolor}[1]{\textcolor{ObjectColor}{\textbf{#1}}}

% --- TikZ styles for ray tracing visualizations ---
\tikzstyle{process} = [rectangle, rounded corners=3mm, minimum width=2cm, minimum height=0.8cm, text centered, draw=PrimaryColor, thick, fill=PrimaryColor!15, drop shadow]
\tikzstyle{arrow} = [thick, PrimaryColor, ->, >=stealth]
\tikzstyle{ray} = [thick, RayColor, ->, >=stealth]
\tikzstyle{lightray} = [thick, LightColor, ->, >=stealth]
\tikzstyle{reflectray} = [thick, SecondaryColor, ->, >=stealth]
\tikzstyle{refractray} = [thick, AccentColor, ->, >=stealth]
\tikzstyle{shadowray} = [thick, DarkGray, ->, >=stealth, dashed]

% --- 3D object styles ---
\tikzstyle{sphere} = [circle, minimum size=1.5cm, draw=ObjectColor, thick, fill=ObjectColor!20, drop shadow]
\tikzstyle{plane} = [rectangle, minimum width=3cm, minimum height=0.2cm, draw=ObjectColor, thick, fill=ObjectColor!20]
\tikzstyle{triangle} = [regular polygon, regular polygon sides=3, minimum size=1.5cm, draw=ObjectColor, thick, fill=ObjectColor!20]

% --- Eye/Camera styles ---
\tikzstyle{eye} = [circle, minimum size=0.8cm, draw=PrimaryColor, thick, fill=PrimaryColor!30]
\tikzstyle{pixel} = [rectangle, minimum size=0.2cm, draw=AccentColor, fill=AccentColor!30]

% --- Text box styles ---
\tikzstyle{conceptbox} = [rectangle, rounded corners, fill=PrimaryColor!10, draw=PrimaryColor, thick, text width=0.8\textwidth, inner sep=8pt]
\tikzstyle{formulabox} = [rectangle, rounded corners, fill=SecondaryColor!10, draw=SecondaryColor, thick, inner sep=6pt]

% --- Custom environments ---
\newtcolorbox{raybox}[1]{
  colback=RayColor!10,
  colframe=RayColor,
  title=#1,
  fonttitle=\bfseries,
  sharp corners
}

\newtcolorbox{conceptbox}[1]{
  colback=PrimaryColor!10,
  colframe=PrimaryColor,
  title=#1,
  fonttitle=\bfseries,
  rounded corners
}

\newtcolorbox{mathbox}[1]{
  colback=SecondaryColor!10,
  colframe=SecondaryColor,
  title=#1,
  fonttitle=\bfseries,
  rounded corners
}

\newenvironment{timeline}{%
  \begin{tikzpicture}[scale=0.8]
    \coordinate (start) at (0,0);
    \newcounter{timelineitem}
    \setcounter{timelineitem}{0}
  }{%
  \end{tikzpicture}
}

\newcommand{\timelineitem}[4]{%
  \stepcounter{timelineitem}%
  % Name of this item’s node
  \edef\thisID{time\thetimelineitem}%
  \ifnum\value{timelineitem}=1
  % First item: anchor at (start)
  \node[rectangle, draw, fill=PrimaryColor!20,
  text width=1.5cm, minimum height=0.8cm, align=center]
  (\thisID) at (start) {\textbf{#2}};
  \else
  % Later items: below=#1 of previous time node
  \pgfmathtruncatemacro\prev{\value{timelineitem}-1}%
  \node[rectangle, draw, fill=PrimaryColor!20,
    text width=1.5cm, minimum height=0.8cm, align=center,
  below=#1 of time\prev]
  (\thisID) {\textbf{#2}};
  \fi
  % Description always to the right of this time node
  \node[rectangle, draw, fill=SecondaryColor!10,
    text width=4cm, minimum height=0.8cm,
  align=left, right=0.3cm of \thisID]
  (desc\thetimelineitem) {\textbf{#3}\\#4};
  % Draw arrow back to previous if not the first
  \ifnum\value{timelineitem}>1
  \draw[->, thick, PrimaryColor]
  (time\prev.south) -- (\thisID.north);
  \fi
}

\tikzset{
  camera/.style={fill=PrimaryColor!60, draw=PrimaryColor!80, rectangle, minimum size=8pt},
  image plane/.style={fill=AccentColor!10, draw=AccentColor!50, opacity=0.8},
  pixel/.style={fill=AccentColor!60, thick},
  primary ray/.style={->, very thick, red!90},
  object/.style={fill=ObjectColor!60, draw=ObjectColor!80, circle, minimum size=12pt},
  fovangle/.style={<->, thick, PrimaryColor, dashed}
}

\tikzset{
  lens/.style={thick, PrimaryColor, line width=3pt},
  focal plane/.style={thick, AccentColor},
  object ray/.style={->, thick, ObjectColor},
  image ray/.style={->, thick, SecondaryColor},
  optical axis/.style={dashed, gray}
}

\newcommand{\lens}[3]{%
  \begingroup
  % half‐dimensions
  \pgfmathsetlengthmacro{\a}{#2}%
  \pgfmathsetlengthmacro{\b}{#3}%
  \pgfmathsetlengthmacro{\c}{#3/2}%
  \begin{scope}[shift={#1}]
    \draw[line join=round, fill=blue!15]
    (0,-{\c})
    arc(-30:30:{\a} and {\b})
    arc(150:210:{\a} and {\b})
    ;
  \end{scope}
  \endgroup
}

% ================================================================
% Reference theme, rebuilt to match combined/deck.pdf
% Rust title bar, cream ground, Adventor headings, thin footer.
% Placed last so it overrides metropolis.
% ================================================================
\usepackage{tgheros}
\usepackage{tgadventor}
\renewcommand{\sfdefault}{qhv}                       % Heros for body copy
\newcommand{\headfont}{\fontfamily{qag}\selectfont}  % Adventor for headings

\definecolor{ThemeRust}{HTML}{9B4632}
\definecolor{ThemeCream}{HTML}{F8F6F2}
\definecolor{ThemeSage}{HTML}{C6CCB1}
\definecolor{ThemeStone}{HTML}{E7E7DB}
\definecolor{ThemeMuted}{HTML}{8B8A88}
\definecolor{ThemeInk}{HTML}{191919}

\setbeamercolor{background canvas}{bg=ThemeCream}
\setbeamercolor{normal text}{fg=ThemeInk}
\setbeamercolor{frametitle}{bg=ThemeRust, fg=ThemeCream}
\setbeamercolor{footer}{bg=ThemeCream, fg=ThemeMuted}
\setbeamercolor{alerted text}{fg=ThemeRust}
\setbeamercolor{itemize item}{fg=ThemeRust}
\setbeamercolor{itemize subitem}{fg=ThemeRust}
\setbeamercolor{palette primary}{bg=ThemeRust, fg=ThemeCream}  % section pages

% full-bleed rust bar across the top
\setbeamertemplate{frametitle}{%
  \nointerlineskip%
  \begin{beamercolorbox}[wd=\paperwidth,sep=0pt]{frametitle}%
    \vbox{%
      \vskip0.60em
      \hbox{\hskip0.9em%
        \begin{minipage}{\dimexpr\paperwidth-1.8em\relax}%
          \raggedright\headfont\bfseries\large\insertframetitle%
        \end{minipage}}%
      \vskip0.55em
    }%
  \end{beamercolorbox}%
}

% thin three-zone footer
\setbeamertemplate{footline}{%
  \hbox{%
  \begin{beamercolorbox}[wd=.40\paperwidth,ht=1.9ex,dp=1.0ex,left]{footer}%
    \hspace*{0.9em}\tiny\insertshortauthor
  \end{beamercolorbox}%
  \begin{beamercolorbox}[wd=.28\paperwidth,ht=1.9ex,dp=1.0ex,center]{footer}%
    \tiny\insertshorttitle
  \end{beamercolorbox}%
  \begin{beamercolorbox}[wd=.32\paperwidth,ht=1.9ex,dp=1.0ex,right]{footer}%
    \tiny\insertframenumber\,/\,\inserttotalframenumber\hspace*{0.9em}
  \end{beamercolorbox}}%
  \vskip1pt
}

% full-bleed rust title page
\setbeamertemplate{title page}{%
  \begin{tikzpicture}[remember picture, overlay]
    \fill[ThemeRust] (current page.south west) rectangle (current page.north east);
  \end{tikzpicture}%
  \vfill
  \begin{flushright}
    {\headfont\bfseries\LARGE\color{ThemeCream}\inserttitle\par}
    \vspace{0.3em}
    {\color{ThemeCream!85}\small\insertsubtitle\par}
    \vspace{1.3em}
    {\color{ThemeCream}\small\insertauthor\par}
    \vspace{0.5em}
    {\color{ThemeCream!80}\scriptsize\insertinstitute\par}
    \vspace{0.8em}
    {\color{ThemeCream!75}\scriptsize\insertdate\par}
  \end{flushright}
  \vfill
}

% --- Deck palette remapped onto the theme -------------------------------
% Overrides the navy/violet set defined at the top of this file. Every TikZ
% figure and tcolorbox references these names, so they all follow along.
\definecolor{PrimaryColor}{HTML}{9B4632}     % rust: highlights, frames, titles
\definecolor{SecondaryColor}{HTML}{7A5C4B}   % warm brown: secondary boxes
\definecolor{AccentColor}{HTML}{7E8C57}      % sage: arrows, alerted text
\definecolor{ObjectColor}{HTML}{A9745B}      % clay: solid shapes
\definecolor{LightColor}{HTML}{D9A441}       % warm gold
\definecolor{TextColor}{HTML}{191919}
\definecolor{BackgroundColor}{HTML}{F8F6F2}
\definecolor{LightGray}{HTML}{E7E7DB}
\definecolor{DarkGray}{HTML}{6B6257}

% the four fern rules -- must match COLS in part 3/figs5.py
\definecolor{RuleA}{HTML}{C0392B}
\definecolor{RuleB}{HTML}{2E7D32}
\definecolor{RuleC}{HTML}{B8860B}
\definecolor{RuleD}{HTML}{1F6F8B}



\title{Fractals}
\subtitle{Infinite Details from Simple Rules}
\author{Speaker 3: Fractals in the Real World}   % <-- put the speaker's name here
\institute{Department of Computer Science and Engineering\\
  Bangladesh University of Engineering and Technology (BUET)}
\date{CSE 200 Presentation}

% one-line takeaway box used on every slide
\newcommand{\takeaway}[1]{%
  \begin{tcolorbox}[colback=PrimaryColor!8, colframe=PrimaryColor, arc=2mm,
    boxrule=0.8pt, left=6pt, right=6pt, top=3pt, bottom=3pt]
    \small #1
  \end{tcolorbox}}
% small muted source line
\newcommand{\source}[1]{{\scriptsize\color{DarkGray}#1}}

\begin{document}

\begin{frame}
  \titlepage
\end{frame}

% ================================================================
% PART 3 — FRACTALS IN THE REAL WORLD (Speaker 3)
% 9 frames: section title, nature, coastline (2), body, database,
%   game worlds, limits, conclusion  (+ references)
% ================================================================
\section{Fractals in the Real World}

% ---------------------------------------------------------------
% S3 Frame 1: Transition — section title (same style as Speaker 2)
% ---------------------------------------------------------------
\begin{frame}{}
  \vfill
  \begin{center}
    {\large\color{SecondaryColor}\itshape
      We know what a fractal is, and how to measure it\ldots}

    \vspace{0.8cm}
    \begin{tikzpicture}
      \draw[PrimaryColor, line width=1.5pt] (-3.5,0) -- (3.5,0);
      \foreach \x in {-3,-2.5,...,3}
        \draw[AccentColor, line width=0.8pt] (\x,0.12) -- (\x,-0.12);
    \end{tikzpicture}

    \vspace{0.8cm}
    {\fontsize{26}{32}\selectfont\bfseries\color{PrimaryColor}
      Fractals in the Real World}

    \vspace{0.5cm}
    {\small\color{DarkGray}
      Where does all of this actually \highlight{show up}?}
  \end{center}
  \vfill
\end{frame}

% ---------------------------------------------------------------
% S3 Frame 2: Nature — real photos next to a rule-grown tree
% ---------------------------------------------------------------
\begin{frame}{Nature draws with one rule, repeated}
  \begin{columns}[T]
    \begin{column}{0.62\textwidth}
      \centering
      \includegraphics[width=0.42\textwidth]{images/part3/nature_romanesco.png}\hfill
      \includegraphics[width=0.42\textwidth]{images/part3/nature_snowflake.png}\\[3pt]
      \includegraphics[width=0.42\textwidth]{images/part3/nature_aloe.png}\hfill
      \includegraphics[width=0.42\textwidth]{images/part3/nature_pinecone.png}
    \end{column}
    \begin{column}{0.36\textwidth}
      \centering
      
\includegraphics[width=0.75\textwidth]{images/tree/7.png}\\[2pt]
      {\footnotesize\color{DarkGray} grown by a computer from\\
        one rule: $F \rightarrow T[-F][+F]$}
    \end{column}
  \end{columns}
  \vspace{0.2cm}
  \takeaway{Broccoli, ice, leaves, cones: the \highlight{same shape at every zoom} ---
    exactly what the Sierpinski rule did on paper.}
\end{frame}

% ---------------------------------------------------------------
% S3 Frame 3: Coastline — Speaker 1's question, answered with data
% ---------------------------------------------------------------
\begin{frame}{A finer ruler makes Britain's coast longer}
  \centering
  \includegraphics[width=\textwidth]{images/part3/britain_rulers.png}

  \vspace{0.1cm}
  \takeaway{This is Speaker~1's question, on \highlight{real map data}:
    the finer the ruler, the more bays and inlets count.}

  \source{Coastline data: GSHHG (Wessel \& Smith 1996), plotted at three resolutions.}
\end{frame}

% ---------------------------------------------------------------
% S3 Frame 4: Richardson plot — Speaker 2's dimension, from data
% ---------------------------------------------------------------
\begin{frame}{Plot it on log--log axes: the slope is the dimension}
  \begin{columns}[c]
    \begin{column}{0.68\textwidth}
      \centering
      \includegraphics[width=\textwidth]{images/part3/richardson.png}
    \end{column}
    \begin{column}{0.31\textwidth}
      \begin{center}
        {\fontsize{26}{30}\selectfont\bfseries\color{PrimaryColor} $D \approx 1.3$}\\[6pt]
        {\small\color{DarkGray} Mandelbrot (1967)\\ found $1.25$}
      \end{center}
      \vspace{0.3cm}
      \takeaway{Speaker~2's box-counting idea, done with a ruler instead of boxes.}
    \end{column}
  \end{columns}
  \vspace{0.2cm}
  \source{Ruler-walk on the full-resolution GSHHG outline of Great Britain; rulers 4--256 km.}
\end{frame}

% ---------------------------------------------------------------
% S3 Frame 5: The body — branching packs area into volume
% ---------------------------------------------------------------
\begin{frame}{Your body branches to pack a huge surface into a small space}
  \begin{columns}[c]
    \begin{column}{0.58\textwidth}
      \centering
      \includegraphics[width=0.62\textwidth]{images/part3/retina.png}\hfill
      \raisebox{0.14\textwidth}{\includegraphics[width=0.34\textwidth]{images/part3/retina_zoom.png}}\\[2pt]
      {\footnotesize\color{DarkGray} blood vessels in a human eye, and a zoom-in}
    \end{column}
    \begin{column}{0.42\textwidth}
      \small
      \begin{itemize}
        \item Lungs: about \highlight{23 splits}\\
          $\Rightarrow$ $\sim$\highlight{70 m$^2$} of surface\\
          inside $\sim$5 litres
        \item Vessels: every cell within reach, using little space
        \item Doctors track the \highlight{dimension} of eye vessels to spot disease early
      \end{itemize}
    \end{column}
  \end{columns}
  \vspace{0.05cm}
  \takeaway{Branching again and again is nature's cheapest way to get \highlight{area out of volume}.}

  \source{Photo: M.~Häggström, ``Fundus photograph of normal left eye'', 2014, CC0. Lung figures: Weibel (1963).}
\end{frame}

% ---------------------------------------------------------------
% S3 Frame 6: Databases — space-filling curve = spatial index
% ---------------------------------------------------------------
\begin{frame}{A curve that fills a square turns a map into a sorted list}
  \begin{columns}[c]
    \begin{column}{0.5\textwidth}
      \centering
      \includegraphics[width=0.95\textwidth]{images/part3/hilbert_index.png}
    \end{column}
    \begin{column}{0.48\textwidth}
      \small
      \begin{itemize}
        \item A database can only sort by \highlight{one number}
        \item Walk the \highlight{Hilbert curve} through the map: each cell gets its position on the curve
        \item Nearby places $\rightarrow$ nearby numbers $\rightarrow$ one quick range search
      \end{itemize}
      \vspace{0.2cm}
      \takeaway{This trick powers map search in systems like Google~S2 and Uber~H3.}
    \end{column}
  \end{columns}
\end{frame}

% ---------------------------------------------------------------
% S3 Frame 7: Game worlds — grown from a rule, not stored
% ---------------------------------------------------------------
\begin{frame}{Game worlds are grown from a rule, not stored}
  \centering
  \includegraphics[width=0.98\textwidth]{images/part3/terrain_octaves.png}

  \vspace{0.15cm}
  \takeaway{Add \highlight{smaller bumps at every scale} --- the Koch idea --- and a
    handful of numbers becomes an endless mountain range. Store the rule, not the map.}

  \source{Terrain generated here by summing 7 layers of noise, each half the size of the last (fractal Brownian motion).}
\end{frame}

% ---------------------------------------------------------------
% S3 Frame 8: Limits — real fractals stop
% ---------------------------------------------------------------
\begin{frame}{Real fractals stop at some scale}
  \begin{columns}[c]
    \begin{column}{0.66\textwidth}
      \centering
      \includegraphics[width=\textwidth]{images/part3/richardson_limit.png}
    \end{column}
    \begin{column}{0.33\textwidth}
      \small
      \begin{itemize}
        \item Coast: bays look like bays, but \highlight{sand grains don't}
        \item Lungs: stop after $\sim$23 splits
        \item Only the Sierpinski triangle goes on \highlight{forever}
      \end{itemize}
      \vspace{0.2cm}
      \takeaway{A fractal is a \highlight{model} of nature --- and every model has a range where it works.}
    \end{column}
  \end{columns}
\end{frame}

% ---------------------------------------------------------------
% S3 Frame 9: Conclusion — the whole talk in one row
% ---------------------------------------------------------------
\begin{frame}{One rule, repeated, explains what smooth geometry couldn't}
  \begin{columns}[T]
    \begin{column}{0.33\textwidth}
      \centering
      {\small\bfseries\color{SecondaryColor} Part 1 --- the surprise}\\[4pt]
      \includegraphics[width=0.7\textwidth]{images/sierpinski/5.png}\\[4pt]
      {\small zero area,\\ infinite edge}
    \end{column}
    \begin{column}{0.33\textwidth}
      \centering
      {\small\bfseries\color{SecondaryColor} Part 2 --- the number}\\[10pt]
      {\fontsize{22}{26}\selectfont\color{PrimaryColor}
        $D=\dfrac{\log N}{\log(1/r)}$}\\[12pt]
      {\small Sierpinski $\approx 1.585$\\ Koch $\approx 1.262$}
    \end{column}
    \begin{column}{0.33\textwidth}
      \centering
      {\small\bfseries\color{SecondaryColor} Part 3 --- the world}\\[4pt]
      \includegraphics[width=0.45\textwidth]{images/britain.png}\\[4pt]
      {\small coast, lungs, maps,\\ game worlds}
    \end{column}
  \end{columns}

  \vspace{0.15cm}
  \begin{center}
    \begin{tikzpicture}
      \node[conceptbox, text width=0.85\textwidth] {
        \centering
        \textit{``Clouds are not spheres, mountains are not cones, coastlines are not
        circles\ldots''} --- Mandelbrot \\[3pt]
        \highlight{Roughness has a number, and simple rules make it.}
      };
    \end{tikzpicture}
  \end{center}
  \begin{center}{\small\color{DarkGray}\itshape Thank you --- questions welcome.}\end{center}
\end{frame}

% ---------------------------------------------------------------
% References
% ---------------------------------------------------------------
\begin{frame}{References}
  \footnotesize
  \begin{itemize}
    \item B.~Mandelbrot, ``How long is the coast of Britain? Statistical self-similarity and fractional dimension,'' \textit{Science}, 1967.
    \item B.~Mandelbrot, \textit{The Fractal Geometry of Nature}, 1982.
    \item P.~Wessel and W.~H.~F.~Smith, ``A global, self-consistent, hierarchical, high-resolution shoreline database,'' \textit{J.~Geophys.~Res.}, 1996 (GSHHG data).
    \item E.~R.~Weibel, \textit{Morphometry of the Human Lung}, Springer, 1963.
    \item M.~Häggström, ``Medical gallery of Mikael Häggström 2014,'' \textit{WikiJournal of Medicine}, 2014 (retina photograph, CC0).
    \item D.~Hilbert, ``\"Uber die stetige Abbildung einer Linie auf ein Fl\"achenst\"uck,'' \textit{Math.~Ann.}, 1891.
    \item A.~Fournier, D.~Fussell, L.~Carpenter, ``Computer rendering of stochastic models,'' \textit{Comm.~ACM}, 1982 (fractal terrain).
    \item 3Blue1Brown, ``Fractals are typically not self-similar'' (video).
  \end{itemize}
\end{frame}

\end{document}
